\begin{tikzpicture}[x=1bp,y=-1bp,draw=black,text=black]
\begin{scope}
\clip (67.333,447.333) rectangle (553.333,729.167);
\q{105}{674}{708}{12.5}{① 下面是小夏同学星期一至星期五上学路上所花时间的统计表。}
\begin{scope}[scale=.6666667]
\fill[white] (129,710) rectangle (271,807);
\draw[black,line width=.5pt] (129,710)--(129,807);
\draw[black,line width=.5pt] (271,710)--(271,807);
\draw[black,line width=.5pt] (381,710)--(381,807);
\draw[black,line width=.5pt] (488,710)--(488,807);
\draw[black,line width=.5pt] (590,710)--(590,807);
\draw[black,line width=.5pt] (693,710)--(693,807);
\draw[black,line width=.5pt] (795,710)--(795,807);
\draw[black,line width=.5pt] (129,710)--(795,710);
\draw[black,line width=.5pt] (129,759)--(795,759);
\draw[black,line width=.5pt] (129,807)--(795,807);
\node[align=center,inner sep=0pt,font=\normalsize] at (200,734.5) {星期};
\node[align=center,inner sep=0pt,font=\normalsize] at (326,734.5) {一};
\node[align=center,inner sep=0pt,font=\normalsize] at (434.5,734.5) {二};
\node[align=center,inner sep=0pt,font=\normalsize] at (539,734.5) {三};
\node[align=center,inner sep=0pt,font=\normalsize] at (641.5,734.5) {四};
\node[align=center,inner sep=0pt,font=\normalsize] at (744,734.5) {五};
\node[align=center,inner sep=0pt,font=\normalsize] at (200,783) {时间/分};
\node[align=center,inner sep=0pt,font=\normalsize] at (326,783) {15};
\node[align=center,inner sep=0pt,font=\normalsize] at (434.5,783) {17};
\node[align=center,inner sep=0pt,font=\normalsize] at (539,783) {14};
\node[align=center,inner sep=0pt,font=\normalsize] at (641.5,783) {16};
\node[align=center,inner sep=0pt,font=\normalsize] at (744,783) {18};
\end{scope}
\q{132}{826}{674}{12.5}{小夏平均每天上学路上要花多长时间？}
\q{104}{984}{723}{12.5}{② 琳琳参加歌唱比赛，8位评委打出的分数分别是94分、98分、95分、}
\q{104}{1024}{723}{12.5}{　95分、88分、96分、93分、97分。去掉一个最高分和一个最低分，其}
\q{104}{1064}{723}{12.5}{　他分数的平均数就是最后得分。琳琳的最后得分是多少？}
\end{scope}
\pgfresetboundingbox\path[use as bounding box] (67.333,447.333) rectangle (553.333,729.167);
\end{tikzpicture}
